\documentclass[10pt,a4paper]{amsart}
\usepackage[margin=19mm]{geometry}
\usepackage{amsmath,amssymb,mathtools}
\usepackage[scr=rsfs,cal=euler]{mathalfa}
\usepackage{xfrac}
\usepackage[unicode,hidelinks,bookmarks=false]{hyperref}
\newcommand{\ie}{i.e.\ }
\newcommand{\cf}{cf.\ }
\newcommand{\resp}{respectively\ }
\newcommand{\Subsection}{\subsection}
\providecommand{\nobreakdash}{\nobreak\hbox{-}}
\setlength{\emergencystretch}{2em}
\multlinegap0pt
\usepackage[shortlabels]{enumitem}
\usepackage{mathtools}
\usepackage{amssymb}
\usepackage{tikz}
\newtheorem{lemma}{Lemma}
\newtheorem{cor}{Corollary}
\newtheorem{thm}{Theorem}
\DeclareMathOperator{\Supp}{Supp}
\newcommand{\g}{\mathfrak{g}}
\title{Weyl modules for Equivariant map Lie superalgebras}
\author{Lakshmi S K, Saudamini Nayak}
\date{Source: arXiv:2511.01631v2 (CC BY 4.0)}
\begin{document}
\maketitle

\noindent\emph{Source and licence.} Weyl modules for Equivariant map Lie superalgebras. Version arXiv:2511.01631v2 (CC BY 4.0), distributed by arXiv under the Creative Commons Attribution 4.0 International licence, http://creativecommons.org/licenses/by/4.0/, as recorded in the arXiv licence metadata for this exact version and shown on the abstract page https://arxiv.org/abs/2511.01631v2. Full licence notice: \texttt{licenses/arxiv-2511.01631-CC-BY-4.0.txt}.\par\smallskip

\noindent\textit{Editorial scope.} Counted unit: the article's thm thm:3 (source0.tex lines 1276--1286). The article's complete original proof of it is delivered with it (source0.tex lines 1287--1321, one \texttt{proof} environment of 35 lines). No mathematics is written, repaired or re-derived: the statement and the proof are the article's own lines, in the article's own notation, and no retained line is substituted or re-flowed.  Retained uncounted as the source-local context the proof needs: lines 1119--1124; lines 1133--1137; lines 1215--1217.  Nothing was omitted from any retained span, so the package carries no omission record.  No retained line contains a citation, so the wrapper carries no bibliography and no bibliography entry.  No retained line contains a figure environment, an image-inclusion command or drawing code, so no image is delivered and the package declares no asset dependency.  Editorial formatting only, in addition: an AMS \texttt{amsart} preamble that restates the article's own declarations and macros used by the retained lines, namely the environments \texttt{lemma}, \texttt{cor}, \texttt{thm} on separate counters, together with the standard packages those lines need.  The retained environments print in this document as Lemma 1, Corollary 1, Theorem 1 in order of appearance, whereas the article prints the same items with the numbering of its own full document.  The complete native source of the article as distributed in the arXiv e-print for this exact version is bundled unchanged as \texttt{sources/188342/source0.tex}.
\begin{lemma}\label{lem:10}
    If $\eta:X_{rat}\rightarrow\mathbb{N}$ satisfies $\Supp~\eta\in X_{*}$, then
    \begin{equation*}
        (\g\otimes I_{\eta})^{\Gamma}=(\g\otimes \tilde I_{\eta})^{\Gamma}, \quad \mbox{where}\quad \tilde I_{\eta}=\prod_{x\in \Supp\eta}\prod_{\gamma\in\Gamma}\textbf{m}_{\gamma.x}^{\eta(x)}.
    \end{equation*}
\end{lemma}

\begin{cor}\label{cor:9.8}
    \begin{equation*}
        (\g\otimes A)^{\Gamma}/(\g\otimes I_{\eta})^{\Gamma}\xrightarrow{\cong}(\g\otimes A)/(\g\otimes I_{\eta}).
    \end{equation*}
\end{cor}

\begin{equation}\label{eq:5}
    (\g\otimes A)\rightarrow(\g\otimes A)/(\g\otimes I_{\eta})\cong(\g\otimes A)^{\Gamma}/(\g\otimes I_{\eta})^{\Gamma}\rightarrow\mbox{End}(V)
\end{equation}

\begin{thm}\label{thm:3}
    For $\mathbf{x}\in X_{*}$, the twisting and the untwisting functors have the following properties.
    \begin{enumerate}[label=(\alph*)]
        \item The twisting functor $\textbf{T}_{\mathbf{x}}$ maps the isomorphism class $\mbox{ev}_{\Psi}$ for $\Psi\in\varepsilon$, $\Supp\Psi\in X_{*}$, to the isomorphism class $\mbox{ev}^{\Gamma}_{\Psi^{\Gamma}}$ for  $\Psi^{\Gamma}\in\varepsilon^{\Gamma}$, where
        \begin{equation*}
            \Psi^{\Gamma}(x)=\sum_{\gamma\in\Gamma}\gamma \Psi(\gamma^{-1}x),\quad x\in X_{rat}.
        \end{equation*}
        \item The untwisting functor $\textbf{U}_{\mathbf{x}}$ maps the isomorphism class of $\mbox{ev}^{\Gamma}_{\Psi}$, $\Psi\in\varepsilon^{\Gamma}$, to the isomorphism class $\mbox{ev}_{\Psi_{\mathbf{x}}}$.
        \item The functor $\textbf{T}_{\mathbf{x}}$ and $\textbf{U}_{\mathbf{x}}$ are mutually inverse isomorphism of categories.
    \end{enumerate}
\end{thm}

\begin{proof}
\begin{enumerate}[label=(\alph*)]
    \item Let $\Psi\in\varepsilon$ and $\Supp\ \Psi\in X_{*}$ and take $\Psi^{\Gamma}$ as 
        $\Psi^{\Gamma}(x)=\sum_{\gamma\in\Gamma}\gamma\Psi(\gamma^{-1}x).$
    For some $\mu\in\Gamma$,
\begin{equation*}
        \mu\Psi^{\Gamma}(x)=\sum_{\gamma\in\Gamma}\mu\gamma\Psi(\gamma^{-1}x)
   \quad \mbox{and} \quad
\Psi^{\Gamma}(\mu x)=\sum_{\gamma\in\Gamma}\gamma \Psi(\gamma^{-1}(\mu x)).
    \end{equation*}
    Let $\mu^{-1}\gamma=\tau$ then $\tau^{-1}=\gamma^{-1}\mu$. Clearly
    \begin{align*}
         \Psi^{\Gamma}(\mu x)&=\sum_{\gamma\in\Gamma}\mu(\mu^{-1}\gamma)\Psi((\mu^{-1}\gamma)^{-1} x)=\sum_{\tau\in\Gamma}\mu\tau\Psi(\tau^{-1}x)\\
        &=\mu\sum_{\tau\in\Gamma}\tau\Psi(\tau^{-1}x)=\mu\Psi^{\Gamma}(x),
    \end{align*}
    and hence $\Psi^{\Gamma}(x)\in\varepsilon^{\Gamma}$. From the definition of $\Psi^{\Gamma}$ we can see that  $\Supp(\Psi^{\Gamma})\subseteq\Gamma.\Supp(\Psi)$. Let $\textbf{x}\in X_{*}$ and $V\in\mathcal{F}_{\textbf{x}}$ is irreducible and corresponds to $\Psi\in\varepsilon$. Let $\mbox{ev}_{\Psi}=(\otimes_{x\in\textbf{x}}\rho_x)\circ\mbox{ev}_{\textbf{x}}$ be the corresponding evaluation representation. Then we have seen that this representation factors through $(\g\otimes A)/(\g\otimes I_{\textbf{x}})$. Using Corollary \ref{cor:9.8} we get $\textbf{T}_{\textbf{x}}(V)$ to be the $(\g\otimes A)^{\Gamma}$-module given by
    \begin{equation}
    \begin{split}
        (\g\otimes A)^{\Gamma}\rightarrow(\g\otimes A)^{\Gamma}/(\g\otimes I_{\textbf{x}})^{\Gamma}\xrightarrow{\cong}(\g\otimes A)/(\g\otimes I_{\textbf{x}})=(\g\otimes A)/(\g\otimes \prod_{x\in\textbf{x}}m_{x})\\\cong\bigoplus_{x\in\textbf{x}}(\g\otimes(A/m_{x}))\xrightarrow{\rho=\otimes_{x\in\textbf{x}}\rho_{x}}End(V).
    \end{split}    
    \end{equation}
    We know already from Lemma \ref{lem:10} that $(\g\otimes I_{\textbf{x}})^{\Gamma}=(\g\otimes \tilde I_{\textbf{x}})^{\Gamma}$ where $\tilde I_{\textbf{x}}=\prod_{x\in\textbf{x}}\prod_{\gamma\in\Gamma}m_{\gamma.x}$. Hence combining this equality with the above equation we get the evaluation representation $(\otimes_{x\in\textbf{x}}\rho_{x})\circ\mbox{ev}^{\Gamma}_{\Gamma.\textbf{x}}$ for $(\g\otimes A)^{\Gamma}$ which is in the isomorphism class of $\mbox{ev}^{\Gamma}_{\Psi^{\Gamma}}$. This completes the proof of (i).
    \item Let $\textbf{x}\in X_{*}$ and $V'\in\mathcal{F}_{\textbf{x}}^{\Gamma}$ be an irreducible module  corresponding to $\Psi\in\varepsilon^{\Gamma}$. Let $\mbox{ev}_{\Psi}^{\Gamma}=(\otimes_{x\in\textbf{x}}\rho_{x})\circ\mbox{ev}_{\textbf{x}}^{\Gamma}$ be the corresponding evaluation representation. Then we get $\textbf{U}_{\textbf{x}}(V')$ to be the $(\g\otimes A)-$ mdoule given by the composition
    \begin{equation*}
        (\g\otimes A)\rightarrow(\g\otimes
        A)/(\g\otimes I_{\textbf{x}})\xrightarrow{\cong}(\g\otimes A)^{\Gamma}/(\g\otimes I_{\textbf{x}})^{\Gamma}\cong\g^{\textbf{x}}\xrightarrow{\otimes_{x\in\textbf{x}}\rho_{x}}End(V').
    \end{equation*}
    This is precisely the evaluation representation of $(\g\otimes A)$ and lies in the isomorphism class of $\mbox{ev}_{\Psi_\textbf{x}}$. Hence (ii) follows.
    \item Suppose $V\in \mathcal{F}_{\textbf{x}}$. Then $V$ is annihilated by some $(\g\otimes I_{\eta})$ and the action of $(\g\otimes A)$ on $\textbf{U}_{\textbf{x}}\textbf{T}_{\textbf{x}}(V)$ is given by
    \begin{equation*}
        (\g\otimes A)\rightarrow(\g\otimes A)/(\g\otimes I_{\eta})\xrightarrow{\cong}(\g\otimes A)^{\Gamma}/(\g\otimes I_{\eta})^{\Gamma}\xrightarrow{\cong}(\g\otimes A)/(\g\otimes I_{\eta})\rightarrow End(V)
    \end{equation*}
    where the two isomorphisms are mutually inverse. Hence $\textbf{U}_{\textbf{x}}\textbf{T}_{\textbf{x}}(V)=V$. Let $\beta$ be a morphism in $\mathcal{F}_{\textbf{x}}$. So $\beta:V\rightarrow W$ be a homomorphism of $(\g\otimes A)-$modules. There exists an $\eta: X_{rat}\rightarrow\mathbb{N}$ with support contained in $\textbf{x}$ such that $(\g\otimes I_{\eta})$ annihilates both $V$ and $W$. By restricting to $(\g\otimes A)^{\Gamma}$, we get $\beta$ to be a morphism in $\mathcal{F}_{\textbf{x}}^{\Gamma}$ denoted by $\textbf{T}_{\textbf{x}}(\beta)$. If we denote $V'$ and $W'$ to be the corresponding $(\g\otimes A)^{\Gamma}$-modules, then again this action of $(\g\otimes A)^\Gamma$ factors through $(\g\otimes A)^{\Gamma}/(\g\otimes I_{\eta})^{\Gamma}\cong(\g\otimes A)/(\g\otimes I_{\eta})$, giving us the morphism $\textbf{T}_{\textbf{x}}(\beta)$ to be a morphism from $V'_{\eta}$ to $W'_{\eta}$. This is exactly $\textbf{U}_{\textbf{x}}\textbf{T}_{\textbf{x}}(\beta)$, morphism of $(\g\otimes A)-$modules. From \eqref{eq:5} it is clear that $V'_{\eta}=V$ and $W'_{\eta}=W$. Hence $\textbf{U}_{\textbf{x}}\textbf{T}_{\textbf{x}}$ is an identity on morphisms and is therefore the identity functor on $\mathcal{F}_{\textbf{x}}$. Similarly $\textbf{T}_{\textbf{x}}\textbf{U}_{\textbf{x}}$ is the identity functor on $\mathcal{F}^{\Gamma}_{\textbf{x}}$. Hence the proof.   
\end{enumerate}
\end{proof}

\end{document}
